\section{Calculus}
\label{sec:calculus}

Algebra computes with expressions as static objects; calculus computes with the
\emph{change} in them. This is the longest section of the book, because it is
where a computer algebra system does its most characteristic work---turning the
limit processes of the calculus into finite, exact algebra. We climb in the
order the subject is usually taught, but with the machinery on show throughout:
differentiation, limits, power series, residues, then the two faces of
integration---first the indefinite integral and the deep decision problem
beneath it, then the definite integral and the complex-analytic and
transform methods that evaluate it---and finally differential equations, where
all of the preceding tools are put to work at once.

The theme that runs through everything is a single distinction. \emph{Some}
operations of calculus are algorithmic in the strict sense: there is a procedure
that always terminates with the right answer or a proof that none exists.
Differentiation is the obvious one---the chain rule is a finite recursion on the
structure of an expression---but so, remarkably, is integration in elementary
terms, by the Risch algorithm. \emph{Other} operations are not decidable at all,
and there the system is a toolbox of specialised engines, each complete on the
family it recognises and honest enough to decline the rest. You will see both
kinds, and the book will be explicit about which is which.

\subsection{Differentiation}
\label{ssec:calc-deriv}

Differentiation\index{differentiation} is the one operation in this section that
is purely mechanical, and Mathilda treats it accordingly. \B{D} computes partial
derivatives, holding every symbol other than the differentiation variable
constant.

\begin{codepairs}
\pair{calculus/derivatives/1}{\B{D}\texttt{[f, x]} differentiates \(f\) with
respect to \(x\): the power rule, \(3x^2\).}
\pair{calculus/derivatives/2}{The chain rule for a composite, \(\frac{d}{dx}\sin
x^2 = 2x\cos x^2\).}
\pair{calculus/derivatives/3}{\texttt{D[f, \{x, n\}]} takes the \(n\)-th
derivative---here the third, \(-a^3\cos ax\).}
\pair{calculus/derivatives/4}{The base-\(b\) logarithm: \(\frac{d}{dx}\log_b x =
1/(x\log b)\).}
\pair{calculus/derivatives/5}{For an \emph{unknown} function the chain rule
produces a \B{Derivative} object: \(f'(g(x))\,g'(x)\), written
\texttt{Derivative[1][\ldots]}.}
\pair{calculus/derivatives/6}{\B{Derivative}\texttt{[2][f]} is the second-derivative
operator itself, carried symbolically until \(f\) is known.}
\end{codepairs}

The last two lines expose the internal representation. An \(n\)-th derivative of
an unknown function is not left as \(D[\ldots]\); it is stored as the operator
\texttt{Derivative[n][f]} applied to its argument, so that
\texttt{Derivative[1][f][g[x]]} means ``\(f'\) evaluated at \(g(x)\).'' This is
what lets Mathilda differentiate the same expression again and again and always
recognise what it is looking at.

Handed a list of variables, \B{D} produces the objects of vector calculus
directly. A single bracketed variable list is a gradient; a repeated one is a
higher array derivative; a vector argument is differentiated component-wise.

\begin{codepairs}
\pair{calculus/derivatives/7}{\texttt{D[f, \{\{x, y\}\}]} is the
gradient\index{gradient} \(\nabla f = (f_x, f_y)\).}
\pair{calculus/derivatives/8}{Repeating the variable list gives the second array
derivative---the Hessian\index{Hessian matrix} \(\nabla^2 f\).}
\pair{calculus/derivatives/9}{For a \emph{vector} field the same call yields the
Jacobian\index{Jacobian matrix}, one gradient row per component.}
\pair{calculus/derivatives/10}{\B{Dt} is the \emph{total} derivative: every
symbol is an implicit function of everything, so each contributes a
\texttt{Dt} differential.}
\pair{calculus/derivatives/11}{So \texttt{Dt} of \(x^n\) carries a \(\log\) term
from the exponent's own dependence.}
\end{codepairs}

The difference between \B{D} and \B{Dt}\index{total derivative} is exactly the
difference between a partial and a total derivative in an ordinary calculus
course: \texttt{D[y, x]} is \(0\) because \(y\) is held fixed, while
\texttt{Dt[y, x]} is the differential \texttt{Dt[y, x]} because \(y\) may depend
on \(x\). Give a function a definition, and its derivative computes to an
ordinary expression:

\begin{codepairs}
\pair{calculus/derivatives/13}{Once \(f(x)=x^5+6x^3\) is defined,
\texttt{f'[x]} evaluates the operator: \(18x^2+5x^4\).}
\pair{calculus/derivatives/14}{And \texttt{f'[5]} is a plain number.}
\end{codepairs}

\calloutnote{Under the hood}{Differentiation is the textbook example of the
system's two-tier design---C for speed, rewrite rules for mathematics---and of
that design evolving. The derivative rules once lived as pattern-matching
\texttt{DownValues} in \texttt{src/internal/deriv.m}; every call walked a linear
list of some sixty rules through the full rule engine. They were promoted to a
native C dispatcher (\texttt{src/calculus/deriv.c}) that branches on the head
symbol in a single comparison and builds the derivative tree directly, for a
2--8\(\times\) speedup, and \texttt{deriv.m} is now an empty hook you can still
drop custom \texttt{Dt} identities into. The higher tier survives where it earns
its keep: the integral tables (\texttt{CRCMathTablesIntegrals.m}) and the mixed-tower
integrator (\texttt{src/internal/mixed/ParallelMixed.m}) below are still Mathilda
programs, loaded at startup.}

The vector operators \B{Grad}, \B{Div}, \B{Curl}, and \B{Laplacian} are assembled
entirely from \B{D}, so they inherit every rule it knows.\index{vector calculus}

\begin{codepairs}
\pair{calculus/vector-calculus/1}{\B{Grad} of a scalar is its gradient vector.}
\pair{calculus/vector-calculus/2}{\B{Grad} of a vector field is its Jacobian.}
\pair{calculus/vector-calculus/3}{\B{Div} contracts a field against the variables:
\(\nabla\cdot\mathbf{F}\).}
\pair{calculus/vector-calculus/4}{\B{Curl} of \((-y, x, 0)\) is the constant
vorticity \((0,0,2)\).}
\pair{calculus/vector-calculus/5}{\B{Laplacian} of \(x^2+y^2+z^2\) is \(6\).}
\pair{calculus/vector-calculus/6}{With a chart name the operators use the
orthonormal basis---here the Coulomb field \(-\nabla(kq/r)= (kq/r^2)\,\hat r\) in
spherical coordinates.}
\pair{calculus/vector-calculus/7}{And the spherical divergence of \(r^2\hat r\)
is \(4r\), from \(\frac{1}{r^2}\partial_r(r^2\cdot r^2)\).}
\end{codepairs}

\calloutnote{Theory}{The three-argument chart form does the real differential
geometry. In a curvilinear orthogonal chart the operators are not just
partial derivatives: each carries the chart's scale factors (Lam\'e
coefficients) \(h_i\), so the spherical divergence is
\(\frac1{J}\sum_i \partial_i\!\big(\tfrac{J}{h_i}F_i\big)\) with \(J=\prod_i
h_i\). Mathilda knows the Cartesian, polar, cylindrical, and spherical charts;
field ranks that would need Christoffel symbols (a vector Laplacian, the gradient
of a vector field in a curved chart) are left unevaluated rather than guessed.}

\subsection{Limits}
\label{ssec:calc-limit}

\B{Limit} evaluates the limit of an expression as a variable approaches a point,
finite or infinite. The easy cases are the ones a first course drills; the hard
ones are where a CAS earns its name.

\begin{codepairs}
\pair{calculus/limits/1}{The archetypal \(0/0\) form: \(\sin x / x \to 1\).}
\pair{calculus/limits/2}{A removable singularity---the difference quotient of
\(x^2\) at \(1\).}
\pair{calculus/limits/3}{Another \(0/0\), giving \(\tfrac12\).}
\pair{calculus/limits/4}{The definition of \(e\), as a \(1^\infty\) form.}
\pair{calculus/limits/5}{With no direction, \(1/x\) at \(0\) is
\mcode{ComplexInfinity}---the two sides disagree in sign.}
\pair{calculus/limits/6}{Fixing \texttt{Direction} takes the one-sided limit,
here \(+\infty\) from above.}
\end{codepairs}

The default is the two-sided limit; \texttt{Direction -> "FromAbove"} and
\texttt{"FromBelow"} select a one-sided approach, and the answer
\mcode{ComplexInfinity} in pair~5 is Mathilda telling you the unsigned limit is
infinite because the signed one does not exist. So far this is bookkeeping.
The following are not:

\begin{codepairs}
\pair{calculus/limits/7}{A \(\infty^0\) form whose value is the largest base,
\(5\)---no series expansion reaches this.}
\pair{calculus/limits/8}{A tower of nested exponentials whose leading terms
cancel to leave \(-1\).}
\pair{calculus/limits/9}{An oscillation that genuinely has no limit:
\mcode{Indeterminate}, not a guessed value.}
\end{codepairs}

Pairs~7 and~8 are the province of a genuine algorithm.\calloutnote{Theory}{The
last resort for a limit of an exp--log function is \emph{Gruntz's most-rapidly-varying
algorithm}\index{Gruntz's algorithm} (his 1996 ETH thesis, implemented in
\texttt{src/calculus/gruntz.c}). Rather than expand from the bottom up---where the
intermediate expressions swell catastrophically on cancellation-heavy towers like
pair~8---it identifies the single most rapidly varying subexpression \(w\), rewrites
the whole function as a series in \(w\to 0^+\), and reads the limit off the leading
term. Because the rewriting is driven by the comparability classes of the
subexpressions, the swell never happens: \((3^x+5^x)^{1/x}\) collapses at once to
\(5\), the base that dominates. Gruntz also runs as the silent fallback inside the
default cascade, after the substitution, rational, series, and L'Hospital layers
have had their turn.} Pair~9 is the honest counterpart: \(x\sin x\) oscillates
with growing amplitude, so no limit exists, and Mathilda proves this (via an
oscillatory normal-form analysis) rather than returning a plausible-looking
number. That refusal to guess is the through-line of the whole section.

Computing a limit is one thing; \emph{proving} one from its definition is another,
and Mathilda does both. The statement \(\lim_{x\to a} f(x) = L\) unfolds into the
\(\varepsilon\)--\(\delta\) sentence\index{limit!epsilon-delta proof of@\(\varepsilon\)--\(\delta\) proof of}\index{epsilon-delta@\(\varepsilon\)--\(\delta\) definition}
\[
  \forall\,\varepsilon>0 \quad \exists\,\delta>0 \quad \forall x \quad
  \bigl(0<\lvert x-a\rvert<\delta \;\Rightarrow\; \lvert f(x)-L\rvert<\varepsilon\bigr),
\]
a quantified formula over the reals---and \B{Resolve} decides quantified formulas
over the reals. Transcribed verbatim, with \B{ForAll} and \B{Exists} carrying each
bound variable and its restriction just as the quantifiers do on the page, the
definition of \(\lim_{x\to2}(3x-1)=5\) is proved outright---and the same sentence
with the limit misstated is \emph{disproved}, not merely left unconfirmed:

\begin{codepairs}
\pair{calculus/limits/10}{The \(\varepsilon\)--\(\delta\) definition of
\(\lim_{x\to2}(3x-1)=5\), transcribed literally and returned \mcode{True}: a proof.}
\pair{calculus/limits/11}{The identical sentence with \(5\) replaced by the wrong
value \(6\)---\mcode{False}. The decision is two-sided.}
\end{codepairs}

Behind the verdict is quantifier elimination, not sampling.\calloutnote{Under the
hood}{An \(\varepsilon\)--\(\delta\) statement is three alternating quantifiers over
the reals, and settling it is \emph{quantifier
elimination}\index{quantifier elimination}: \B{Resolve} eliminates the innermost
block to a condition on the variables that remain, then the next, until only
\mcode{True} or \mcode{False} is left. The engine underneath
(\texttt{src/solve/reduce\_cad.c}) is a \emph{cylindrical algebraic
decomposition}\index{cylindrical algebraic decomposition}---a complete decision
procedure for the first-order theory of real-closed
fields\index{real-closed field} (Tarski's theorem)---so for any
\emph{semialgebraic} \(f\) (polynomial, rational, or algebraic) the answer is
guaranteed, and each \(\lvert\cdot\rvert\) is absorbed by splitting the absolute
value into its two sign branches. That completeness has a sharp boundary: a
transcendental target such as \(\lim_{x\to0}\sin(x)/x = 1\) has no polynomial
decomposition, so Mathilda \emph{declines}---returns the sentence
unevaluated---rather than guess. \B{Limit} and \B{Resolve} are complementary:
\B{Limit} reaches the transcendental limits a decision procedure cannot, while
\B{Resolve} certifies the algebraic ones with a proof.}

\usagebox{Limit}

\subsection{Series}
\label{ssec:calc-series}

A power series is the local linearisation carried to all orders. \B{Series}
computes the Taylor, Laurent, or Puiseux expansion of an expression about a
point, to a requested order.

\begin{codepairs}
\pair{calculus/series/1}{The Taylor series\index{Taylor series} of \(e^x\) to
order five, with its \(O[x]^6\) remainder term.}
\pair{calculus/series/2}{\B{Normal} strips the \(O\) term, leaving an ordinary
polynomial you can compute with.}
\pair{calculus/series/3}{\B{SeriesCoefficient} extracts a single coefficient
without forming the whole series---here \(1/7!\).}
\pair{calculus/series/4}{A \emph{Laurent}\index{Laurent series} series: \(1/\sin
x\) has a pole at \(0\), so the expansion starts at \(x^{-1}\).}
\pair{calculus/series/5}{A \emph{Puiseux}\index{Puiseux series} series in
fractional powers, for \(\sqrt{\sin x}\).}
\pair{calculus/series/6}{Symbolic exponents survive: the binomial series of
\((1+x)^n\) keeps \(n\) unevaluated.}
\end{codepairs}

The three expansion types are one mechanism seen at three scales: a Taylor series
when the function is analytic, a Laurent series with negative powers at a pole, a
Puiseux series with fractional powers at a branch point. The expansion can also be
taken at infinity, and an unknown function expands formally through its
derivatives:

\begin{codepairs}
\pair{calculus/series/7}{An asymptotic\index{asymptotic expansion} expansion at
\(\infty\): the arctangent tends to \(\pi/2\) with a \(1/x\) tail.}
\pair{calculus/series/8}{For an unknown \(f\) the coefficients appear as
\B{Derivative} values---the Taylor series in symbolic form.}
\pair{calculus/series/9}{The internal object is \B{SeriesData}; constructing one
by hand prints as the same series.}
\end{codepairs}

\calloutnote{Under the hood}{Behind the printed sum is a \B{SeriesData} object
carrying the expansion variable, the centre, the coefficient list, the leading and
trailing exponents, and a denominator that makes fractional powers exact integers
internally. It is a first-class value: \B{D}, \B{Integrate}, \B{Plus}, \B{Times},
and \B{Power} all operate on it term by term and return a new truncated series, so
you can do arithmetic with series as freely as with polynomials. Coefficients
promote automatically to arbitrary-precision rationals when machine integers would
overflow, so a deep Laurent expansion such as \(\operatorname{Series}[1/\sin^{10}x]\)
stays exact.}

\subsection{Residues}
\label{ssec:calc-residue}

The residue of a function at an isolated singularity---the coefficient of
\((z-z_0)^{-1}\) in its Laurent series---is the single number that complex
analysis needs to integrate around a pole. \B{Residue} reads it straight out of
the series.

\begin{codepairs}
\pair{calculus/residues/1}{The residue\index{residue} of \(1/z\) at \(0\) is
\(1\).}
\pair{calculus/residues/2}{A double pole with no simple part: \(1/z^2\) has
residue \(0\).}
\pair{calculus/residues/3}{\(\cot z\) has a simple pole at \(0\) with residue
\(1\).}
\pair{calculus/residues/4}{A higher-order pole: residue \(-3/4\) at the double
pole \(z=0\).}
\pair{calculus/residues/5}{A complex pole---\(1/(z^2+1)\) at \(z=i\) has residue
\(-i/2\).}
\pair{calculus/residues/6}{A fifth-order pole: \(1/\sin^5 z\) at \(0\)
gives \(3/8\).}
\pair{calculus/residues/7}{At a \emph{branch} point the residue is undefined, so
the call is left unevaluated.}
\end{codepairs}

\calloutnote{Under the hood}{For a general integrand the residue is the \(-1\)
coefficient of \B{Series}; the engine raises the expansion order automatically
when the pole's order is hidden behind an unknown factor. But a rational function
with an algebraic (radical) pole location would trip the series inverter, so a
simple pole short-circuits to the exact formula \(P(z_0)/Q'(z_0)\) computed
directly, and higher-order rational poles expand the numerator and denominator
\emph{separately} to keep each algebraic coefficient in its minimal radical
spelling. A fractional-power (Puiseux) expansion signals a branch point---pair~7---where
the notion of residue does not apply, so the call declines rather than return a
Laurent coefficient that is not there.} The residue theorem this feeds---summing
residues to evaluate a real integral---reappears in \S\ref{ssec:calc-definite},
where Mathilda uses it as a definite-integration engine in its own right.

\subsection{Indefinite integration: the Risch story}
\label{ssec:calc-integrate}

Integration is the centre of this section and the deepest single algorithm in the
system. Where differentiation is a finite recursion, integration in closed form
is a genuine decision problem, and the theory that settles it is one of the
triumphs of twentieth-century computer algebra. We start with the cases that look
familiar and work toward the machinery that makes them possible.

\begin{codepairs}
\pair{calculus/integrate-elementary/1}{Term-by-term integration of a polynomial.}
\pair{calculus/integrate-elementary/2}{A logarithmic derivative recognised
directly: \(\int 2x/(x^2+1) = \log(1+x^2)\).}
\pair{calculus/integrate-elementary/3}{The arctangent, from an irreducible
quadratic denominator.}
\pair{calculus/integrate-elementary/4}{A partial-fraction case; the answer combines
two of the three logarithms into a single quadratic \(\log\).}
\pair{calculus/integrate-elementary/5}{Two irreducible quadratic factors give a
mixed \(\log\)/\(\arctan\) answer.}
\end{codepairs}

These are all \emph{rational} integrands, and for them integration is completely
solved.\calloutnote{Under the hood}{The rational integrator
(\texttt{src/calculus/intrat.c}) is a faithful implementation of the modern
pipeline: \emph{Hermite reduction}\index{Hermite reduction} splits off the part
of the answer that is itself rational (the repeated-pole part), leaving a
residual with a square-free denominator; the \emph{Lazard--Rioboo--Trager}
algorithm\index{Rothstein--Trager} then produces the logarithmic part as a sum
\(\sum_i c_i\log(g_i)\) whose constants \(c_i\) are the roots of a
resultant---never by an undetermined-coefficient solve. Rioboo's conversion turns
each complex-conjugate pair of logarithms into a real \(\arctan\), which is why
pair~5 comes out in real form. For plain rational inputs the heavy polynomial
arithmetic is handed to FLINT~\cite{flint}; the classical account is
Geddes--Czapor--Labahn~\cite{gcl}. Every branch ends at the same universal check:
\(\operatorname{Cancel}[\operatorname{Together}[D[\text{result}]-f]]\) must be
\(0\).} Rational functions of radicals are next, rationalised by the right
substitution:

\begin{codepairs}
\pair{calculus/integrate-elementary/6}{A rational function of \(\sqrt x\),
rationalised by \(u=\sqrt x\).}
\pair{calculus/integrate-elementary/7}{A square root of a quadratic, handled by a
Euler substitution.}
\end{codepairs}

Now the mathematics deepens. Many perfectly innocent-looking integrands have
\emph{no} antiderivative among the elementary functions---and this is a theorem,
not a gap. \emph{Liouville's theorem}\index{Liouville's theorem} says that if an
elementary function has an elementary antiderivative, that antiderivative has a
very restricted shape: the original function plus a sum of constant multiples of
logarithms of elementary functions. When no such shape exists, the integral is
provably non-elementary, and the honest answer is a \emph{new} function named for
exactly that integral.

\begin{codepairs}
\pair{calculus/integrate-special/1}{Integration by parts, done for you:
\(\int x e^x = (x-1)e^x\).}
\pair{calculus/integrate-special/2}{\(e^{-x^2}\) has no elementary integral, so
the answer is the error function \B{Erf}.}
\pair{calculus/integrate-special/3}{\(e^x/x\) gives the exponential integral
\B{ExpIntegralEi}.}
\pair{calculus/integrate-special/4}{\(1/\log x\) is the logarithmic integral,
returned as \(\operatorname{Ei}(\log x)\).}
\pair{calculus/integrate-special/5}{\(\sin x/x\) gives the sine integral
\B{SinIntegral}.}
\pair{calculus/integrate-special/6}{A dilogarithm\index{dilogarithm} appears:
\(\int \log x/(1+x) = \log x\log(1+x) + \operatorname{Li}_2(-x)\).}
\pair{calculus/integrate-special/7}{And a genuine polylogarithm ladder for
\(x/(e^x-1)\), the Bose--Einstein integrand.}
\end{codepairs}

Each of these is correct \emph{by construction}: the algorithm certifies that no
elementary form exists before it reaches for the special function, so \B{Erf},
\B{ExpIntegralEi}, \B{SinIntegral}, and \B{PolyLog} are not approximations or
placeholders---they are the exact antiderivatives, and their appearance is a
statement about the integrand.\calloutnote{Theory}{The engine behind pairs~2--7
is the \emph{Risch algorithm}\index{Risch algorithm}
(\texttt{src/calculus/integrate\_risch\_transcendental.c}), the decision
procedure for integration in a differential field built by adjoining
exponentials and logarithms one at a time. Its heart is the
\emph{Risch differential equation}: closing the exponential part of an integral
amounts to solving \(q' + f'q = p\) for a polynomial \(q\), which Mathilda does by
Bronstein's polynomial-time SPDE reduction rather than an undetermined-coefficient
blow-up. Layered above the base algorithm are engines for the special-function
answers: a Gaussian recogniser for \B{Erf}, Cherry's algorithms for the exponential
and logarithmic integrals and the \emph{polylogarithm ladder}\index{polylogarithm
ladder} \(\int x^n/(e^{cx}-\rho)\), and a dilogarithm recogniser. When the field
decision proves an integral non-elementary but the special-function engines cannot
name it, the integral is returned unevaluated with an \texttt{Integrate::nonelem}
message explaining why.}

The book promised to be explicit about what is decidable, and integration lets it
keep the promise directly. The predicate that decides elementary integrability is
exposed:

\begin{codepairs}
\pair{calculus/elementary-q/1}{\(e^x/x\) is \emph{not} elementary---a proof, not a
timeout.}
\pair{calculus/elementary-q/2}{Nor is \(e^{x^2}\).}
\pair{calculus/elementary-q/4}{But \(e^x/(1+e^x)\) is: its integral is
\(\log(1+e^x)\).}
\pair{calculus/elementary-q/5}{And \(\sin x/x\) is not---the sine integral is the
statement that it is not.}
\end{codepairs}

A \texttt{False} here is backed by an exact certificate---a non-constant residue,
or a Risch differential equation with no rational solution---so it means ``no
elementary antiderivative exists,'' never ``I gave up.'' This soundness is
purchased at a real cost, which is worth seeing:

\begin{codepairs}
\pair{calculus/elementary-q/6}{A correct answer in an awkward dress:
\(\int\cos(\log x)\) is real, but printed through complex powers of \(x\).}
\end{codepairs}

\calloutnote{Pitfall}{The output of pair~6 is exactly
\(\tfrac{x}{2}(\cos\log x + \sin\log x)\)---but Mathilda returns it as
\((\tfrac14-\tfrac{i}4)x^{1+i}+(\tfrac14+\tfrac{i}4)x^{1-i}\), a genuinely real
quantity written with conjugate complex exponents. The transcendental integrator
reaches trigonometric integrands by exponentialising with \(\cos =
(e^{i\cdot}+e^{-i\cdot})/2\), and the current simplifier does not always fold the
result back to real form. The value is right and \(\B{N}\) confirms it; the form is
a known cosmetic limitation, not an error. Prefer the algebraic routes when a real
closed form matters.}

Two further capabilities of \B{Integrate} deserve mention, because they reach
integrands the classical Risch algorithm alone cannot. The first is the
\texttt{"ParallelMixedTower"}\index{ParallelMixedTower method} method, a parallel
Risch--Norman integrator over a \emph{mixed} transcendental tower that may contain
a square root anywhere in its structure. It is what closes hard algebraic-function
integrals such as the classic \(\int\sqrt{\tan x}\):

\begin{codepairs}
\pair{calculus/integrate-special/8}{\(\int\sqrt{\tan x}\,dx\): a mixed radical
tower, closed in real \(\arctan\)/\(\log\) form.}
\end{codepairs}

The parallel method treats the whole integrand as living in a differential field
built over rational operations, one square root, and \(\log\)/\(\exp\)/\(\tan\) of
field elements, and solves for the antiderivative by an ansatz whose unknowns are
found by linear algebra---then verifies every result by re-differentiating before
returning it, so a case it cannot realise correctly declines rather than escaping
as a wrong answer. The second is Feynman's trick,
\texttt{"DiffUnderInt"}---differentiation under the integral sign---which belongs
to definite integration and is taken up next.

\usagebox{Integrate}

\subsection{Definite integration}
\label{ssec:calc-definite}

A definite integral is not just an antiderivative evaluated at two points---not
when the integrand has poles, the interval runs to infinity, or the
antiderivative crosses a branch cut. Mathilda's definite integrator is a cascade
of methods, each correct on the family it claims. The workhorse is the
fundamental theorem of calculus, made careful.

\begin{codepairs}
\pair{calculus/definite/1}{The elementary case: \(\int_0^1 x^2 = 1/3\).}
\pair{calculus/definite/2}{An \emph{improper} integral that converges---the
singularity at \(0\) is integrable.}
\pair{calculus/definite/3}{An infinite upper limit, evaluated through the
\B{Limit} engine.}
\pair{calculus/definite/11}{An iterated double integral, reduced innermost-first.}
\end{codepairs}

\calloutnote{Under the hood}{The Newton--Leibniz mechanism
(\texttt{src/calculus/integrate\_newton\_leibniz.c}) antidifferentiates \(f\),
locates the real poles of the \emph{integrand} strictly inside the interval, splits
the interval there, and evaluates each boundary through \B{Limit}---a plain limit at
infinite endpoints, a one-sided limit at interior poles. This is how the improper
integrals above acquire their finite values, and how a genuinely divergent
integral is caught and left unevaluated rather than assigned the difference of two
infinities.} That last safeguard matters:

\begin{codepairs}
\pair{calculus/definite/10}{\(\int_{-1}^1 dx/x\) diverges, and Mathilda says so by
returning the input unevaluated (with an \texttt{Integrate::idiv} message).}
\end{codepairs}

For the classical improper and periodic integrals that complex analysis dispatches
by summing residues, a residue-theorem engine runs \emph{before} Newton--Leibniz.
These answers are correct by construction---once the structural gates hold, the
residue theorem gives the exact value.

\begin{codepairs}
\pair{calculus/definite/4}{A rational function over the whole line, by residues in
the upper half-plane: \(\pi/\sqrt2\).}
\pair{calculus/definite/5}{A Fourier integral, by Jordan's lemma\index{Jordan's
lemma}: \(\int_{-\infty}^\infty \cos x/(1+x^2) = \pi/e\).}
\pair{calculus/definite/9}{A rational function of \(\cos x\) over a full period, via
\(z=e^{ix}\) on the unit circle.}
\end{codepairs}

A third family---half-line integrals of transcendental functions that neither
residues nor the FTC close---is handled by a Mellin-transform engine built on the
Ramanujan Master Theorem\index{Ramanujan Master Theorem}, and a fourth by
symmetry.

\begin{codepairs}
\pair{calculus/definite/6}{The Gaussian \(\int_0^\infty e^{-x^2} = \sqrt\pi/2\).}
\pair{calculus/definite/7}{The sinc integral \(\int_0^\infty \sin x/x = \pi/2\).}
\pair{calculus/definite/8}{The Debye integral \(\int_0^\infty x^3/(e^x-1) =
\pi^4/15\)---the Stefan--Boltzmann law, from the Bose--Einstein kernel.}
\end{codepairs}

\calloutnote{Theory}{Pair~8 is Ramanujan's Master Theorem\index{Mellin transform}
at work. If \(f(x)=\sum_k (-1)^k\varphi(k)x^k/k!\) then its Mellin transform
\(\int_0^\infty x^{s-1}f(x)\,dx\) equals \(\Gamma(s)\varphi(-s)\) on a
fundamental strip, and the Bose--Einstein kernel \(1/(e^{x}-1)\) lands on
\(\Gamma(s)\zeta(s)\)---so \(\int_0^\infty x^3/(e^x-1)=\Gamma(4)\zeta(4)=6\cdot
\pi^4/90=\pi^4/15\). Every application is gated on its convergence strip, checked
by \B{Simplify} or against your \texttt{Assumptions}; when the strip cannot be
proved the value is returned as a \B{ConditionalExpression}, matching Wolfram, and
a provably violated strip declines. The Fermi--Dirac companion
\(\int_0^\infty x^3/(e^x+1)=7\pi^4/120\) comes out the same way.}

When an endpoint is complex, or the spec lists more than two points, \B{Integrate}
switches to a genuine contour integral along the straight segments.

\begin{codepairs}
\pair{calculus/contour/1}{A contour integral of \(1/z\) between two complex
points---a difference of principal logarithms.}
\pair{calculus/contour/2}{\(\int_0^{1+i} z^2\,dz\), path-independent because
\(z^2\) is entire.}
\pair{calculus/contour/3}{A closed square loop around the origin gives \(2\pi
i\)---the residue theorem\index{residue theorem}, numerically.}
\end{codepairs}

The Cauchy principal value\index{Cauchy principal value} is available for the
odd-order interior poles where the symmetric divergences cancel:

\begin{codepairs}
\pair{calculus/contour/4}{The principal value across the pole at \(x=2\):
\(-\log 2\).}
\end{codepairs}

Finally, differentiation under the integral sign---\emph{Feynman's
trick}\index{Feynman's trick}\index{differentiation under the integral sign}---is
a definite-integration method in its own right. For a parameter-dependent integral
\(I(p)\), it differentiates the integrand in \(p\), evaluates the simpler integral,
and integrates back, fixing the constant from an exact base value.

\begin{codepairs}
\pair{calculus/feynman/1}{\(\int_0^1 (x^a-1)/\log x\,dx = \log(1+a)\), the standard
Feynman warm-up.}
\pair{calculus/feynman/2}{And \(\int_0^\infty \log(1+a^2x^2)/(1+x^2)\,dx = \pi\log(1+a)\)
under \(a>0\).}
\end{codepairs}

\subsection{Differential equations}
\label{ssec:calc-dsolve}

The whole apparatus of this section converges on \B{DSolve}, which solves
differential equations symbolically. It returns solutions as rules, with the
constants of integration named \texttt{C[1]}, \texttt{C[2]}, and so on; initial
or boundary conditions, given as equations, are fitted to those constants. Like
\B{Integrate}, it is a cascade polyalgorithm---a long sequence of specialised
methods, each recognising an equation type it can solve completely. We survey it
by that classification, which is also the classification any course on
differential equations teaches.

We begin, as promised, with the warm-up:

\begin{codepairs}
\pair{calculus/dsolve-first-order/1}{The exponential equation \(y'=y\): the
solution is \(C_1 e^x\).}
\pair{calculus/dsolve-first-order/2}{Separable, \(y'=xy\): \(C_1 e^{x^2/2}\).}
\pair{calculus/dsolve-first-order/3}{First-order linear, solved by an integrating
factor\index{integrating factor}.}
\pair{calculus/dsolve-first-order/4}{A Bernoulli\index{Bernoulli equation}
equation \(y'=y+xy^2\), linearised by \(v=y^{-1}\).}
\pair{calculus/dsolve-first-order/5}{An exact\index{exact equation} equation;
the potential inverts to the explicit \(y=C_1/(1+x^2)\).}
\pair{calculus/dsolve-first-order/6}{A Clairaut\index{Clairaut equation} equation,
whose general solution is a family of lines.}
\pair{calculus/dsolve-first-order/7}{A quadratic (Riccati\index{Riccati equation})
right-hand side, whose solution is a shifted tangent.}
\pair{calculus/dsolve-first-order/8}{An initial-value problem---the constant is
fitted to \(y(0)=5\).}
\end{codepairs}

Each line above is a different classical method: separable, linear, Bernoulli,
exact, Clairaut, Riccati. \B{DSolve} recognises the shape of the equation and
dispatches to the matching engine, and every branch it returns is verified by
back-substitution before you see it. The linear constant-coefficient case---the
backbone of a differential-equations course---is next, with the forcing handled by
undetermined coefficients or variation of parameters.

\begin{codepairs}
\pair{calculus/dsolve-linear/1}{A damped oscillator: complex characteristic roots
give \(e^{-2x}\) times sine and cosine.}
\pair{calculus/dsolve-linear/2}{Resonant forcing, \(y''-y=e^x\): the particular
solution carries the tell-tale factor of \(x\).}
\pair{calculus/dsolve-linear/3}{Repeated integration, in \B{Function} form.}
\pair{calculus/dsolve-linear/4}{A boundary-value problem, fitted at two points to
give \(\sin x\).}
\pair{calculus/dsolve-linear/5}{An \emph{inconsistent} BVP: \texttt{\{\}} means no
solution exists, distinct from an unhandled equation.}
\end{codepairs}

The resonance in pair~2 is worth pausing on: when the forcing frequency coincides
with a natural frequency, the naive particular solution fails, and \B{DSolve}
detects the coincidence and multiplies by \(x\) automatically---the same
resonance a physicist watches for in a driven oscillator. Pair~5 shows the
system's care with existence: an over-determined boundary-value problem returns
the empty list as a proof that no solution exists, which is a different statement
from leaving an equation symbolic because no method applied.

Beyond constant coefficients lies the variable-coefficient world, where the
solutions are the named special functions of \S\ref{ssec:calc-series} and the
answer often hinges on a deep classical algorithm.

\begin{codepairs}
\pair{calculus/dsolve-special/1}{The Airy equation \(y''=xy\), whose solutions are
\B{AiryAi} and \B{AiryBi}.}
\pair{calculus/dsolve-special/2}{Bessel's equation of order \(2\), giving
\B{BesselJ} and \B{BesselY}.}
\pair{calculus/dsolve-special/3}{An Euler--Cauchy\index{Euler--Cauchy equation}
equidimensional equation, solved by \(x^r\) with complex exponents.}
\pair{calculus/dsolve-special/4}{A Liouvillian solution found by Kovacic's
algorithm---elementary, not a series.}
\pair{calculus/dsolve-special/5}{No closed form exists, so the fallback is a
Frobenius\index{Frobenius method} power series about the origin.}
\end{codepairs}

\calloutnote{Theory}{Pair~4 is the crown of second-order theory.
\emph{Kovacic's algorithm}\index{Kovacic's algorithm}
(\texttt{src/calculus/dsolve\_kovacic.c}) decides whether a second-order linear
ODE \(y''+Py'+Qy=0\) has a \emph{Liouvillian} solution---one built from exponentials,
integrals, and algebraic functions---and constructs it when it does, by examining
the equation's reduced normal form \(z''=rz\) and the local behaviour of \(r\) at
its poles and at infinity. When it succeeds you get a genuine closed form (here
\((C_1+C_2e^{x^2})e^{-x^2/2}/\sqrt x\)); when it proves none exists, control passes
to the Frobenius\index{Frobenius method} series method, which expands about a regular
singular point with an indicial equation fixing the leading exponents. The two
together---Kovacic for the closed form, Frobenius for its absence---cover the
second-order linear equations a course would set.}

Systems of equations are dispatched to their own cascade, which diagonalises the
coefficient matrix and handles real, complex, and defective spectra.

\begin{codepairs}
\pair{calculus/dsolve-systems/1}{A \(2\times2\) system with complex eigenvalues:
the solution is oscillatory.}
\pair{calculus/dsolve-systems/2}{A system with real eigenvalues \(-1\) and \(3\):
a combination of growing and decaying exponentials.}
\end{codepairs}

Two engines stand out for their generality. The first is the Lie point-symmetry
method, the general backstop for first-order equations.\calloutnote{Theory}{A
\emph{Lie point symmetry}\index{Lie point symmetry} is an infinitesimal
transformation \(V=\xi\,\partial_x+\eta\,\partial_y\) that maps solution curves to
solution curves. Finding one is an underdetermined PDE---so Mathilda's engine
(documented in \texttt{docs/design/dsolve\_lie\_symmetry.md}) runs a table of
ansatz heuristics, each proposing a shape for \((\xi,\eta)\) and testing it against
the linearised symmetry condition. An accepted symmetry hands over an integrating
factor \(\mu=1/(\eta-\xi\omega)\), whose quadrature is the first integral of the
equation. Because a first integral need not admit an explicit \(y(x)\), the answer
is often returned \emph{implicitly}, as a relation \(G(x,y)=C_1\).}

\begin{codepairs}
\pair{calculus/dsolve-lie/1}{A symmetry with a rational infinitesimal yields an
implicit first integral.}
\pair{calculus/dsolve-lie/2}{An irrational right-hand side \(\sqrt{ax+by+c}\),
handled by a similarity-ansatz symmetry.}
\end{codepairs}

The second is the partial-differential-equation solver, which handles first-order
equations by characteristics and the classical second-order equations of
mathematical physics.

\begin{codepairs}
\pair{calculus/dsolve-pde/1}{The transport equation, by the method of
characteristics\index{method of characteristics}: an arbitrary function of
\(x-ct\).}
\pair{calculus/dsolve-pde/2}{A first-order PDE with a lower-order term and forcing.}
\pair{calculus/dsolve-pde/3}{The wave equation as an initial-value problem, by
d'Alembert's formula\index{d'Alembert's formula}.}
\pair{calculus/dsolve-pde/5}{A Sturm--Liouville\index{Sturm--Liouville problem}
eigenvalue problem: the eigenvalues \(n^2\) and eigenfunctions \(\sin(nx)\).}
\end{codepairs}

The arbitrary function \texttt{C[1][\ldots]} in pairs~1 and~2 is the PDE analogue
of an integration constant---a first-order PDE's general solution contains an
arbitrary \emph{function}, not merely a constant, and it is a function of the
characteristic coordinate the equation is constant along. Pair~3 is d'Alembert's
formula for the vibrating string, and pair~5 is the eigenvalue quantisation that
underlies every separated boundary-value problem, returned as a
\B{ConditionalExpression} indexed by an integer. The heat equation is solved by
its Gaussian kernel; because that convolution is non-elementary for a general
initial condition, Mathilda keeps it as an integral rather than inventing a closed
form---the same honesty about non-elementarity that governs \B{Integrate}.

\usagebox{DSolve}

\subsection*{Where this connects}

Calculus is where the rest of Mathilda's mathematics is put to work at once.
Every integral of a rational function leans on the factoring, partial-fraction,
and resultant machinery of \S\ref{sec:algebra}; every \B{Root} object in an
antiderivative is an exact algebraic number the solver returns; the definite
integrals close through the same \B{Limit} engine that opens this section, and the
contour methods rest on \B{Residue}. \B{DSolve} in turn calls on all of it---on
\B{Integrate} for its quadratures, on the eigenvalue routines of \S\ref{sec:linalg}
for its systems, and on the special functions of \S\ref{ssec:calc-series} for its
closed forms. Where a symbolic answer does not exist, or exists but is too costly
to want, the numerical companions \B{NIntegrate}, \B{NDSolve}, and their relatives
in \S\ref{sec:numcalc} take over, and the numeric loops that evaluate a
series or a special function accelerate through the compiler of
Chapter~\ref{ch:compilation}. As always, the exhaustive options for any function
named here are one \texttt{?Name} away at the prompt.
